% 50-11(50쪽) — y축에 대하여 대칭인 이차함수 y=f(x) 의 그래프와 그래프 위의 점 A(2, f(2)) 를 지나는 직선 OA. 스캔 화질이 낮아 TikZ 로 다시 그림.
% 원문에 있는 것만: 포물선(꼭짓점이 원점) · 원점을 지나는 직선 · A 에서 x축으로 내린 점선 · 라벨 y=f(x), A, O, 2, x, y. 눈금·f(2) 값은 원문에 없어 적지 않는다.
\begin{tikzpicture}[x=0.36cm, y=0.40cm, line width=0.5pt]
  \draw[->, >=stealth, line width=0.7pt] (-4.4,0) -- (4.7,0) node[below=1pt] {$x$};
  \draw[->, >=stealth, line width=0.7pt] (0,-1.4) -- (0,4.5) node[left=1pt] {$y$};
  \node[below left=-2pt and -3pt, font=\small] at (0,0) {$\pt{O}$};
  % y축에 대하여 대칭인 이차함수 (꼭짓점 = 원점)
  \draw[line width=0.6pt, domain=-3.72:3.72, samples=90, smooth] plot (\x, {0.25*\x*\x});
  % 직선 OA (A 의 좌우로 연장)
  \draw[line width=0.5pt] (-2.4,-1.2) -- (4.4,2.2);
  % A 에서 x축에 내린 점선
  \draw[densely dashed, line width=0.4pt] (2,1) -- (2,0);
  \node[above left=-1pt and -3pt, font=\small] at (2,1) {$\pt{A}$};
  \node[below=1pt, font=\small] at (2,0) {$2$};
  \node[anchor=west, font=\small] at (2.55,3.95) {$y=f(x)$};
\end{tikzpicture}
